\chapter{Referencial teórico}
\label{cap:referencial_teorico}

Faça uma compilação dos principais autores sobre o tema a ser investigado. Nesta fase é importante o apoio de seu orientador\index{orientador} e a clareza quanto ao tema e à pergunta da sua pesquisa, indicando o que é essencial ser lido e resumido e o que é secundário. A pergunta funciona como um guia do que deve ser exposto nessa seção e, principalmente, do que não é necessário.

A forma pela qual você estruturará esse seção já será uma de suas contribuições. Cada pesquisador, dentro de um mesmo tema, possui uma forma única de organizar o conhecimento compulsado.

O Referencial teórico\index{referencial teórico} é como uma aula sobre o tema para que o leitor consiga acompanhar as próximas seções de seu trabalho. Observe que não é um resumo de tudo que você aprendeu no seu curso.

Divida o assunto que você deseja estudar em seções. As seções devem conduzir o leitor pelos conceitos apresentados. Assim, você deve começar pelo tópico mais geral e ir especificando cada vez mais até trazer seu leitor ao tema específico de seu trabalho acadêmico. Dentro de cada seção, você pode usar subseções para organizar melhor a informação destacada.

\section{Assunto mais geral do referencial teórico} \label{sub:assunto2}

A primeira seção\index{seção} do referencial teórico normalmente será sobre o tema mais geral dentre os assuntos abordados no seu trabalho. 

\section{Assunto mais específico do referencial teórico}

As demais seções tendem a ser mais específicas que as anteriores

\section{Trabalhos relacionados}

Nessa seção, você pode elencar os trabalhos com os quais as suas contribuições podem ser comparadas na seção de discussão ao final do seu documento.
Para que o leitor\index{leitor} entenda as similaridades entre seu trabalho\index{trabalho} e os trabalhos correlatos, você deve, aqui, fazer um resumo de cada trabalho correlato realçando a interseção que eles possuem com os seus objetivos de pesquisa.

\section{Quadro referencial teórico}

Você pode incluir um quadro ao final do seu referencial destacando para o leitor quais foram as suposições ou hipóteses construídas a partir da literatura que você compilou. A palavra hipótese é mais utilizada em pesquisas quantitativas. A palavra suposição é mais utilizada em pesquisas mais interpretativas, qualitativas.

O trabalho \cite{DissertacaoElton2015}, por utilizar uma metodologia mista, apresenta exemplos de quadro de hipóteses e de quadro de suposições. A tabela~\ref{tab:quadro_hipoteses} apresenta uma parte do quadro de hipóteses\index{quadro de hipóteses} do trabalho.

\begin{table}[H]
	\caption{Exemplo de parte de um quadro de hipóteses.}
	\begin{center}
		\fontsize{10pt}{13pt}\selectfont
		\begin{tabular}{|p{0.1\linewidth}|p{0.4\linewidth}|p{0.4\linewidth}|}
			\hline
			\rowcolor{gray!10}
			\begin{center}
				\textbf{Hipótese}
			\end{center} &
			\begin{center}
				\textbf{Declaração}
			\end{center} &
			\begin{center}
				\textbf{Referencial}
			\end{center} \\
			\hline
			{\centering H1 \par} & O Prazer Percebido impacta da Utilidade Percebida. & Dias, 2001; Igbaria, Parasuraman e Baroudi, 1996. \\
			\hline
			{\centering H2 \par} & O Prazer Percebido impacta da Facilidade de Uso Percebida. & Dias, 2001; Igbaria, Parasuraman e Baroudi, 1996. \\
			\hline
			{\centering H3 \par} & O Prazer Percebido impacta na Intenção de Uso. & Dias, 2001; Igbaria, Parasuraman e Baroudi, 1996. \\
			\hline
			\multicolumn{3}{l}{Fonte: Adaptado de \cite[37]{DissertacaoElton2015}}
		\end{tabular}
		\label{tab:quadro_hipoteses}
	\end{center}
\end{table}

Ainda sobre o mesmo trabalho, a figura~\ref{tab:quadro_suposicoes} apresenta uma parte do quadro de suposições\index{quadro de suposições}.

\begin{table}[H]
	\caption{Exemplo de parte de um quadro de suposições.}
	\begin{center}
		\fontsize{10pt}{13pt}\selectfont
		\begin{tabular}{|p{0.1\linewidth}|p{0.4\linewidth}|p{0.4\linewidth}|}
			\hline
			\rowcolor{gray!10}
			\begin{center}
				\textbf{Hipótese}
			\end{center} &
			\begin{center}
				\textbf{Declaração}
			\end{center} &
			\begin{center}
				\textbf{Referencial}
			\end{center} \\
			\hline
			{\centering S1 \par} & Há uma tendência à utilização de celulares, smartphones e tablets como apoio a aprendizagem. & Santos, 2009; Mateus, Brito, 2011; Hoehle, Zhang, Venkatesh, 2015.. \\
			\hline
			{\centering S2 \par} & A presença de tutores presenciais no início do curso para ambientar os novos alunos é importante. & Mateus, Brito, 2011; Bizzaria et al. 2015.. \\
			\hline
			\multicolumn{3}{l}{Fonte: Adaptado de \cite[37]{DissertacaoElton2015}}
		\end{tabular}
		\label{tab:quadro_suposicoes}
	\end{center}
\end{table}
